
\section{Estado del arte}

\label{sec:estado_del_arte}

% TODO: ampliar con trabajos recientes (2019-2026) sobre predicción de delitos con ML
Diversos estudios han aplicado modelos supervisados y no supervisados
para predecir y clasificar la criminalidad en entornos urbanos. Los
métodos de regresión y clasificación clásicos siguen siendo ampliamente
usados por su interpretabilidad \parencite{james2013introduction},
mientras que las redes neuronales profundas han mostrado buenos resultados
cuando se dispone de grandes volúmenes de datos georreferenciados
\parencite{goodfellow2016deep}.

En el contexto peruano existen pocas aplicaciones documentadas de
estos modelos a nivel de distritos urbanos, lo que constituye una
oportunidad para el presente trabajo. (TODO: revisar literatura nacional
y local).
